\documentclass[12pt,a4paper]{article}
\usepackage{amsmath,amssymb,amsthm,hyperref,geometry}
\usepackage{enumitem,booktabs,array}
\usepackage[T1]{fontenc}
\usepackage[utf8]{inputenc}
\geometry{margin=2.5cm}

% -------------------------------------------------------
% Theorem environments
% -------------------------------------------------------
\newtheorem{theorem}{Theorem}[section]
\newtheorem{lemma}[theorem]{Lemma}
\newtheorem{corollary}[theorem]{Corollary}
\newtheorem{proposition}[theorem]{Proposition}
\newtheorem{conjecture}[theorem]{Conjecture}
\newtheorem{definition}[theorem]{Definition}

\theoremstyle{remark}
\newtheorem{remark}[theorem]{Remark}
\newtheorem{example}[theorem]{Example}

% -------------------------------------------------------
% Metadata
% -------------------------------------------------------
\title{\textbf{Navier--Stokes Equations: Existence and Smoothness}\\[0.5em]
\large Weak Solutions, Regularity, Turbulence, and the Millennium Problem}
\author{Michael Fuhrmann}
\date{March 2026}

\hypersetup{
  pdftitle={Navier-Stokes Existence and Smoothness},
  pdfauthor={Michael Fuhrmann},
  colorlinks=true,
  linkcolor=blue,
  citecolor=blue
}

\begin{document}
\maketitle
\begin{abstract}
The Navier--Stokes equations govern the motion of viscous incompressible fluids and
underlie virtually all modelling of fluid dynamics, from aeronautics to weather prediction.
The Millennium Prize Problem asks whether smooth (infinitely differentiable) global
solutions always exist in three dimensions given smooth initial data, or whether
singularities (``blow-up'') can develop in finite time.

This paper presents the equations, the function spaces (Sobolev, Leray) needed for weak
solutions, Leray's fundamental 1934 results, Ladyzhenskaya's complete 2D theory, the
Caffarelli--Kohn--Nirenberg partial regularity theorem, Kolmogorov's 1941 turbulence
theory, and Tao's 2016 averaged Navier--Stokes blow-up construction.  The current
state in 2026 is summarised.
\end{abstract}

\tableofcontents
\newpage

% ================================================================
\section{The Navier--Stokes Equations}
% ================================================================

\subsection{Derivation and Physical Meaning}

The \emph{Navier--Stokes equations} for an incompressible viscous Newtonian fluid with
velocity field $u\colon \mathbb{R}^3 \times [0,\infty) \to \mathbb{R}^3$ and pressure
$p\colon \mathbb{R}^3 \times [0,\infty) \to \mathbb{R}$ are:
\begin{align}
  \partial_t u + (u \cdot \nabla)u &= -\nabla p + \nu \Delta u + f, \label{eq:ns_momentum}\\
  \nabla \cdot u &= 0, \label{eq:ns_incompressible}
\end{align}
with initial condition $u(x, 0) = u_0(x)$.

Here:
\begin{itemize}[nosep]
  \item $(u \cdot \nabla)u$ is the \emph{advection} (nonlinear) term,
  \item $-\nabla p$ is the \emph{pressure gradient},
  \item $\nu > 0$ is the \emph{kinematic viscosity},
  \item $\nu \Delta u$ is the \emph{diffusion} (viscous stress) term,
  \item $f$ is an external body force (e.g.\ gravity).
\end{itemize}

Equation \eqref{eq:ns_incompressible} expresses the incompressibility constraint
$\operatorname{div} u = 0$, which means the fluid is volume-preserving.

\subsection{The Millennium Problem Statement}

The Clay Mathematics Institute~\cite{Clay2000} asks for one of:
\begin{enumerate}[nosep]
  \item A proof that for any smooth, rapidly decaying initial data $u_0 \in C^\infty(\mathbb{R}^3)$
    with $\nabla \cdot u_0 = 0$, there exists a smooth global solution $(u, p)$ with
    $u \in C^\infty(\mathbb{R}^3 \times [0,\infty))$.
  \item A proof of the \emph{existence of a blow-up}: initial data for which no smooth
    global solution exists.
\end{enumerate}

The problem is open in $\mathbb{R}^3$ (and on the torus $\mathbb{T}^3$).

% ================================================================
\section{Function Spaces and Weak Formulation}
% ================================================================

\subsection{Sobolev Spaces}

\begin{definition}[Sobolev Space $H^k$]
  For $k \geq 0$, the Sobolev space $H^k(\mathbb{R}^n) = W^{k,2}(\mathbb{R}^n)$ consists
  of functions $f \in L^2(\mathbb{R}^n)$ with weak derivatives up to order $k$ in $L^2$:
  \[
    H^k(\mathbb{R}^n) \;=\; \bigl\{ f \in L^2(\mathbb{R}^n) :
    \|\hat{f}\|_{H^k}^2 := \int_{\mathbb{R}^n} (1+|\xi|^2)^k |\hat{f}(\xi)|^2\,d\xi < \infty \bigr\}.
  \]
\end{definition}

The energy of the Navier--Stokes solution is measured by the $L^2$ norm of $u$:
\[
  E(t) \;=\; \frac{1}{2}\int_{\mathbb{R}^3} |u(x,t)|^2\,dx.
\]
Formally, multiplying \eqref{eq:ns_momentum} by $u$ and integrating yields the
\emph{energy inequality}:
\begin{equation}\label{eq:energy}
  E(t) + \nu \int_0^t \int_{\mathbb{R}^3} |\nabla u|^2\,dx\,ds \;\leq\; E(0),
\end{equation}
for $f \equiv 0$.  This shows energy is non-increasing (viscosity dissipates energy).

\subsection{Leray's Weak Formulation (1934)}

\begin{definition}[Weak (Leray) Solution~\cite{Leray1934}]
  A \emph{weak solution} on $[0,T]$ is a divergence-free vector field
  $u \in L^\infty(0,T; L^2) \cap L^2(0,T; H^1)$ satisfying:
  \begin{enumerate}[nosep]
    \item For all divergence-free test functions $\phi \in C_c^\infty(\mathbb{R}^3 \times (0,T))$:
    \[
      \int_0^T\!\int_{\mathbb{R}^3} \!\bigl(-u \cdot \partial_t \phi
      - (u\otimes u):\nabla\phi + \nu \nabla u : \nabla\phi\bigr)\,dx\,dt = 0.
    \]
    \item The energy inequality \eqref{eq:energy} holds.
    \item $u(t) \to u_0$ in $L^2$ as $t \to 0^+$.
  \end{enumerate}
\end{definition}

\begin{theorem}[Leray 1934~\cite{Leray1934}]
  For any $u_0 \in L^2(\mathbb{R}^3)$ with $\nabla \cdot u_0 = 0$ and $f \in L^2$,
  there exists at least one global weak (Leray) solution.
\end{theorem}

Leray's proof uses a \emph{regularisation} (mollification of the nonlinearity by a
parameter $\varepsilon$), then passes to the limit $\varepsilon \to 0$ using compactness
(Aubin--Lions lemma).  The resulting solution is called a \emph{Leray--Hopf weak solution}
after Leray~\cite{Leray1934} and Hopf~\cite{Hopf1951}.

% ================================================================
\section{The 2D Case: Complete Theory}
% ================================================================

\begin{theorem}[Ladyzhenskaya 1969~\cite{Lady1969}]
  In two space dimensions ($u\colon \mathbb{R}^2 \times [0,\infty) \to \mathbb{R}^2$),
  for any $u_0 \in H^1(\mathbb{R}^2)$ with $\nabla \cdot u_0 = 0$:
  \begin{enumerate}[nosep]
    \item There exists a unique global smooth solution.
    \item The solution depends continuously on initial data.
    \item The flow is globally well-posed in $H^k$ for all $k \geq 1$.
  \end{enumerate}
\end{theorem}

The key ingredient in 2D is the \emph{Ladyzhenskaya inequality}:
\[
  \|u\|_{L^4}^2 \;\leq\; C\|u\|_{L^2}\|\nabla u\|_{L^2},
\]
which, combined with the 2D vorticity equation $\partial_t \omega + (u\cdot\nabla)\omega = \nu\Delta\omega$
(where the vortex-stretching term $\omega\cdot\nabla u$ vanishes in 2D), gives an
$L^\infty$ bound on vorticity and prevents blow-up.

In 3D, the corresponding vorticity equation is
$\partial_t \omega + (u\cdot\nabla)\omega = (\omega\cdot\nabla)u + \nu\Delta\omega$,
where the \emph{vortex-stretching term} $(\omega\cdot\nabla)u$ can amplify vorticity,
potentially leading to blow-up.

% ================================================================
\section{The 3D Regularity Problem}
% ================================================================

\subsection{Local Existence of Smooth Solutions}

\begin{theorem}[Kato--Fujita 1962~\cite{KatoFujita1962}]
  For $u_0 \in H^s(\mathbb{R}^3)$ with $s \geq 1/2$ and $\nabla\cdot u_0=0$, there
  exists $T^* > 0$ and a unique smooth solution $u \in C([0,T^*); H^s)$.
\end{theorem}

The question is whether the maximal existence time $T^*$ is always $+\infty$ (global
existence) or can be finite (blow-up).

\subsection{Blow-Up Criteria}

\begin{theorem}[Beale--Kato--Majda 1984~\cite{BKM1984}]
  If $T^* < \infty$ (blow-up in finite time), then necessarily
  \[
    \int_0^{T^*} \|\omega(\cdot,t)\|_{L^\infty}\,dt = +\infty,
  \]
  where $\omega = \nabla \times u$ is the vorticity.  Equivalently, the vorticity must
  become unbounded for blow-up to occur.
\end{theorem}

This localises the regularity problem to the behaviour of vorticity.

\subsection{Serrin's Regularity Criterion}

\begin{theorem}[Serrin 1962~\cite{Serrin1962}; Prodi 1959]
  Let $u$ be a Leray--Hopf weak solution on $(0,T)$.  If additionally
  $u \in L^q(0,T; L^r(\mathbb{R}^3))$ for some pair $(q,r)$ with
  \[
    \frac{2}{q} + \frac{3}{r} = 1, \quad r > 3,
  \]
  then $u$ is the unique smooth solution on $(0,T)$.
\end{theorem}

The Serrin class bounds are scale-invariant (the Navier--Stokes equations are invariant
under $(u,p,x,t) \mapsto (\lambda u, \lambda^2 p, \lambda^{-1}x, \lambda^{-2}t)$), and
the critical case $r=3$, $q=\infty$ (i.e.\ $u \in L^\infty(0,T; L^3)$) is the deepest
regular endpoint, proved by Escauriaza--Seregin--Šverák (2003)~\cite{ESS2003}.

% ================================================================
\section{Caffarelli--Kohn--Nirenberg Partial Regularity}
% ================================================================

\begin{theorem}[Caffarelli--Kohn--Nirenberg 1982~\cite{CKN1982}]
  Let $(u, p)$ be a \emph{suitable weak solution} on $\mathbb{R}^3 \times (0,T)$ (a
  Leray--Hopf solution satisfying a local energy inequality).  Then the singular set
  $\mathcal{S} = \{(x,t) : u \text{ is not locally bounded near } (x,t)\}$ satisfies
  \[
    \mathcal{H}^1(\mathcal{S}) = 0,
  \]
  where $\mathcal{H}^1$ is the one-dimensional parabolic Hausdorff measure.
  In particular, the singular set has \emph{parabolic Hausdorff dimension} $\leq 1$.
\end{theorem}

This is the strongest unconditional regularity result in 3D.  In particular:
\begin{itemize}[nosep]
  \item The singular set cannot contain any curve, surface, or open set.
  \item At most finitely many singular points can exist at any fixed time $t$.
  \item ``Space-time'' singularities (if any) are very sparse.
\end{itemize}

\begin{remark}
  The CKN theorem does not exclude the possibility that $\mathcal{S}$ is non-empty;
  it merely bounds its size.  Constructing a suitable weak solution with a non-empty
  singular set, or proving $\mathcal{S} = \emptyset$ for all smooth data, remain open.
\end{remark}

% ================================================================
\section{Turbulence: Kolmogorov's 1941 Theory}
% ================================================================

\subsection{The Energy Cascade}

Kolmogorov~\cite{Kolmogorov1941} (1941) proposed a statistical theory of fully developed
turbulence based on two hypotheses:
\begin{enumerate}[nosep]
  \item \textbf{Local isotropy:} At high Reynolds numbers, small-scale turbulent motions
    are statistically isotropic (independent of large-scale geometry).
  \item \textbf{Universal equilibrium:} In the \emph{inertial subrange} (scales between
    the integral scale $L$ and the viscous dissipation scale $\eta$), statistics depend
    only on the mean energy dissipation rate $\varepsilon$ and viscosity $\nu$.
\end{enumerate}

\subsection{The Kolmogorov $-5/3$ Spectrum}

From dimensional analysis alone (using $[\varepsilon] = \text{m}^2\text{s}^{-3}$ and
$[k] = \text{m}^{-1}$), the energy spectrum $E(k)$ (energy per unit wavenumber) must obey:
\begin{equation}\label{eq:k41}
  E(k) \;=\; C_K \varepsilon^{2/3} k^{-5/3}, \quad \eta^{-1} \ll k \ll L^{-1},
\end{equation}
where $C_K \approx 1.5$ is the \emph{Kolmogorov constant}.

The Kolmogorov dissipation scale is $\eta = (\nu^3/\varepsilon)^{1/4}$ and the Reynolds
number $\mathrm{Re} = L^{4/3}\varepsilon^{1/3}/\nu$.  The $k^{-5/3}$ law
\eqref{eq:k41} is confirmed experimentally over many orders of magnitude.

\subsection{K41 and Navier--Stokes}

Rigorous justification of K41 from the Navier--Stokes equations remains elusive.
Constantin and Doering~\cite{ConstantinDoering1994} proved upper bounds on energy
dissipation consistent with K41.  The connection between K41 phenomenology and the
mathematical regularity problem (blow-up) is an active research area.

% ================================================================
\section{Tao's Averaged Navier--Stokes Blow-Up (2016)}
% ================================================================

\subsection{The Averaged Equations}

In 2016, Terence Tao~\cite{Tao2016} constructed a blow-up for a modified system called
the \emph{averaged Navier--Stokes equations}:
\[
  \partial_t u = \nu\Delta u - B(u,u), \quad \nabla\cdot u = 0,
\]
where $B(u,v) = \frac{1}{2}[\nabla\times(u\times v) + P\nabla(u\cdot v)]$ is a
symmetrised bilinear form related to the Navier--Stokes nonlinearity, averaged over
SO(3)-rotations:
\[
  B_{\mathrm{avg}}(u,v) \;=\; \mathbb{E}_{R \in SO(3)}\bigl[R^{-1}B(Ru, Rv)\bigr].
\]

\subsection{Tao's Theorem}

\begin{theorem}[Tao 2016~\cite{Tao2016}]
  There exists an averaged Navier--Stokes system with a smooth, divergence-free initial
  condition $u_0 \in C^\infty(\mathbb{R}^3)$ such that the unique smooth solution
  $u(t)$ blows up at a finite time $T^* < \infty$:
  \[
    \limsup_{t \to T^*} \|u(t)\|_{H^1(\mathbb{R}^3)} = +\infty.
  \]
\end{theorem}

\begin{remark}
  The averaged system satisfies the \emph{energy equality} (not just inequality),
  is divergence-free, and shares the same scaling invariance as Navier--Stokes.
  It differs from the true Navier--Stokes equations in that the nonlinearity is averaged.
\end{remark}

\subsection{Implications of Tao's Result}

Tao's paper~\cite{Tao2016} serves several purposes:
\begin{enumerate}[nosep]
  \item It shows that blow-up is \emph{not ruled out} by the structure of the equations
    alone (energy inequality, divergence-free condition, scaling).
  \item The blow-up mechanism (an ``energy cascade machine'') provides a conceptual model
    for how singularities might form in the true Navier--Stokes system.
  \item It identifies barriers: any proof of global regularity must use some property
    of Navier--Stokes not shared by the averaged system.
\end{enumerate}

% ================================================================
\section{Recent Developments (2020--2026)}
% ================================================================

\begin{itemize}[nosep]
  \item \textbf{Non-uniqueness of weak solutions (2019--2022):} Buckmaster and
    Vicol~\cite{BuckmasterVicol2019} proved using convex integration that non-unique
    weak solutions with $H^{1/3-\varepsilon}$ regularity exist.  This suggests Leray
    solutions are non-unique below a critical regularity threshold.
  \item \textbf{Onsager conjecture (proved):} De Lellis and Székelyhidi (2019)~\cite{DeLellis2019}
    proved the Onsager conjecture: there exist solutions to the \emph{Euler equations}
    ($\nu=0$) with $C^{1/3-\varepsilon}$ regularity that dissipate energy, while all
    solutions with $C^{1/3+\varepsilon}$ conserve energy.
  \item \textbf{Stochastic Navier--Stokes:} Adding noise can restore uniqueness (Da Prato
    et al.); uniqueness for stochastic Navier--Stokes in 3D in some classes was
    established by Hofmanova et al.\ (2022).
  \item \textbf{Machine learning and turbulence (2023--2026):} Neural operators (Fourier
    Neural Operators by Li et al.) have achieved remarkable accuracy in turbulence
    simulation, but these are numerical/statistical approaches unrelated to the mathematical
    existence problem.
  \item \textbf{Millennium Problem (2026):} No proof of global regularity or blow-up for
    the true 3D Navier--Stokes equations has been found.  The problem remains fully open.
\end{itemize}

% ================================================================
\section{Summary of Results}
% ================================================================

\begin{table}[h]
\centering
\begin{tabular}{lll}
\toprule
Dimension & Existence & Uniqueness/Regularity \\
\midrule
2D & Global, smooth (Leray 1934)   & Unique, smooth (Ladyzhenskaya 1969) \\
3D & Global weak (Leray 1934)       & Open \\
3D (short time) & Local smooth (Kato--Fujita 1962) & Unique (short time) \\
3D (singular set) & Hausdorff dim.\ $\leq 1$ (CKN 1982) & --- \\
Averaged 3D & Blow-up (Tao 2016)     & --- \\
\bottomrule
\end{tabular}
\caption{Summary of Navier--Stokes results}
\end{table}

% ================================================================
\section{Conclusion}
% ================================================================

The Navier--Stokes problem is a paradigm of the deep gap between physical intuition
and mathematical rigour.  Fluids behave smoothly in every physical experiment and
numerical simulation, yet proving this behaviour from the equations alone has resisted
all efforts for nearly a century.

The partial results — Leray's weak solutions, Ladyzhenskaya's complete 2D theory,
Caffarelli--Kohn--Nirenberg's singular set bound, and Tao's averaged blow-up — collectively
delineate the boundary of current knowledge.  Resolution of the Millennium Problem will
likely require genuinely new mathematics: new function spaces, new regularity criteria,
or a unexpected blow-up construction.

\begin{thebibliography}{99}

\bibitem{Clay2000}
C.~L.~Fefferman,
\textit{Existence and Smoothness of the Navier--Stokes Equation},
Clay Mathematics Institute Millennium Problem description, 2000,
\url{https://www.claymath.org/millennium-problems/navier-stokes-equation}.

\bibitem{Leray1934}
J.~Leray,
\textit{Sur le mouvement d'un liquide visqueux emplissant l'espace},
Acta Math.\ \textbf{63} (1934), 193--248.

\bibitem{Hopf1951}
E.~Hopf,
\textit{Über die Anfangswertaufgabe für die hydrodynamischen Grundgleichungen},
Math.\ Nachr.\ \textbf{4} (1951), 213--231.

\bibitem{Lady1969}
O.~A.~Ladyzhenskaya,
\textit{The Mathematical Theory of Viscous Incompressible Flow},
Gordon and Breach, New York, 1969.

\bibitem{KatoFujita1962}
T.~Kato and H.~Fujita,
\textit{On the non-stationary Navier--Stokes system},
Rend.\ Sem.\ Mat.\ Univ.\ Padova \textbf{32} (1962), 243--260.

\bibitem{BKM1984}
J.~T.~Beale, T.~Kato, and A.~Majda,
\textit{Remarks on the breakdown of smooth solutions for the 3-D Euler equations},
Comm.\ Math.\ Phys.\ \textbf{94} (1984), 61--66.

\bibitem{Serrin1962}
J.~Serrin,
\textit{On the interior regularity of weak solutions of the Navier--Stokes equations},
Arch.\ Rational Mech.\ Anal.\ \textbf{9} (1962), 187--195.

\bibitem{ESS2003}
L.~Escauriaza, G.~Seregin, and V.~Šverák,
\textit{$L_{3,\infty}$-solutions of the Navier--Stokes equations and backward uniqueness},
Russian Math.\ Surveys \textbf{58} (2003), 211--250.

\bibitem{CKN1982}
L.~Caffarelli, R.~Kohn, and L.~Nirenberg,
\textit{Partial regularity of suitable weak solutions of the Navier--Stokes equations},
Comm.\ Pure Appl.\ Math.\ \textbf{35} (1982), 771--831.

\bibitem{Kolmogorov1941}
A.~N.~Kolmogorov,
\textit{The local structure of turbulence in incompressible viscous fluid for very large Reynolds numbers},
Dokl.\ Akad.\ Nauk SSSR \textbf{30} (1941), 301--305.

\bibitem{ConstantinDoering1994}
P.~Constantin and C.~R.~Doering,
\textit{Variational bounds on energy dissipation in incompressible flows: II. Channel flow},
Phys.\ Rev.\ E \textbf{51} (1995), 3192--3198.

\bibitem{Tao2016}
T.~Tao,
\textit{Finite time blowup for an averaged three-dimensional Navier--Stokes equation},
J.\ Amer.\ Math.\ Soc.\ \textbf{29} (2016), 601--674.

\bibitem{BuckmasterVicol2019}
T.~Buckmaster and V.~Vicol,
\textit{Nonuniqueness of weak solutions to the Navier--Stokes equation},
Ann.\ of Math.\ \textbf{189} (2019), 101--144.

\bibitem{DeLellis2019}
C.~De Lellis and L.~Székelyhidi,
\textit{On turbulence and geometry: from Nash to Onsager},
Notices Amer.\ Math.\ Soc.\ \textbf{66} (2019), 677--685.

\bibitem{ConstantinFoias1988}
P.~Constantin and C.~Foias,
\textit{Navier--Stokes Equations},
University of Chicago Press, 1988.

\bibitem{Temam1984}
R.~Temam,
\textit{Navier--Stokes Equations: Theory and Numerical Analysis},
AMS Chelsea Publishing, 1984.

\end{thebibliography}

\end{document}
